\documentclass{zhvt-classic}

% 使用本地源流明体（GenRyuMin2 TW）按文件路径加载，不依赖系统字体名
% 版框（内框／界栏／书口侧栏／鱼尾）改用与版心印章同一印泥棕色：
%   印章图片 yshs.png 实色 = RGB(135,68,52) = HTML 874434
\usepackage{xcolor}
\usepackage{tikz}
% sealbrown 双注册：xcolor 供 TikZ/pgf，L3 供 cls 版框
\definecolor{sealbrown}{HTML}{874434}
\ExplSyntaxOn
\color_set:nnn { sealbrown } { HTML } { 874434 }
\ExplSyntaxOff

\zhvtset{
  grid color       = sealbrown,
  main font        = GenRyuMin2TW-SB,
  main font path   = C:/Users/Administrator/texmf/fonts/opentype/public/GenRyuMin2TW/,
jiazhu font      = GenRyuMin2TW-M,
  jiazhu font path = C:/Users/Administrator/texmf/fonts/opentype/public/GenRyuMin2TW/,
  % 页码框容纳 4 字：正常页「九十七」+前置空 1 格=4 行恰好；
  % 100+ 页「一百零一」共 5 行超出，末字「一/二/三」被框裁掉（2026-10-04 实测），故放宽到 5
  shukou page number chars = 5,
  % 书口书名框高从默认 6 字减到 4 字，把省出的 2 字高拨给 chapter 框，
  % 使 6 字书名+卷名放进 9 字高的 chapter 框后不再压线（2026-10-05）
  shukou title chars = 4,
}

\title{史记}[外戚世家]
\author{汉·司马迁}
\maker{闲叟漫记·重製}


% ===== 索隐述赞样式（2026-11-08 四次定稿 · 愚大人亲手改定参数）=====
% ⚠ 竖排前提：自定义字体家族必须带 RawFeature={vertical:+vert}；
%   simhei.ttf 加 +vhal 会整体加载失败静默回退，只写 +vert；simfang 同理。
\newCJKfontfamily{\shuzanfont}[Path=C:/Windows/Fonts/,RawFeature={vertical:+vert}]{simhei.ttf}
\newCJKfontfamily{\shuzanbody}[Path=C:/Windows/Fonts/,RawFeature={vertical:+vert}]{simfang.ttf}
% 题头：152x30pt 圆角 6pt 棕底白字；baseline (0,2) 使整框落入界行格子不压版框线；
% 白字节点 x 不补偿、y +8.2pt（像素级墨迹实测 dx/dy≈0，L/R 留白约 20pt）。
\newcommand{\shuzanhead}[1]{%
  \begin{tikzpicture}[baseline={(0,2)}]
    \fill[sealbrown, rounded corners=6pt]
      (-76pt,-15pt) rectangle (76pt,15pt);
    \node[text=white, inner sep=0pt, anchor=center] at (0pt,8.2pt)
      {\shuzanfont\fontsize{23}{23}\selectfont #1};
  \end{tikzpicture}%
}
% ★ 修正书口卷名位置：chapter 框起始位置改回 cls 默认值，避免全局影响。
%   只在每个卷首页通过 \juan 命令里的空命令微调该页位置（安全，不会跨页污染）。
\ExplSyntaxOn
\AtBeginDocument{
  \dim_set:Nn \l__zhvt_shukou_chapter_title_start_dim
    { \l__zhvt_shukou_title_height_dim + \l__zhvt_font_size_dim * 2 }
  \dim_set:Nn \l__zhvt_shukou_page_number_start_dim
    { \l__zhvt_shukou_chapter_title_start_dim + \l__zhvt_shukou_chapter_title_height_dim }
}
\ExplSyntaxOff

\begin{document}

% ★★ 版心（书口）排布覆写：把书名放到鱼尾之下（2026-10-04）
%
%   cls 默认版心自上而下：书名（title 静态框，在鱼尾之上）→ 鱼尾 → 篇名（chapter 动态框）→ 页码
%   本模板改为：          鱼尾（位置不动）→ 书名 → 篇名 → 页码分割线 → 页码
%
%   手法：把 cls 放在 title 静态框里的书名，改放进 chapter 动态框（鱼尾之下）；
%         title 静态框留空。书名取 \l__zhvt_title_str，即 \title{...} 的第一参数，
%         所以换书时只改 \title，不必在此重复写书名。
%
%   ⚠ 三条必须遵守的写法（2026-10-04 实测）：
%     ① 汉字必须用 \hbox{...} 包住。flowfram 会把框内容装进一个宽 14.4mm（= 一个 line_sep
%        = 书口栏宽）的纵向 minipage 并按该宽度折行，而整页靠 PDF /Rotate 90 显示、
%        框内文字本身不变向 —— 不包 \hbox 时「围炉夜话」4 字会被逐字折成 4 行，
%        在版面上就成了横排 4 字并横向溢出版框。
%     ② 竖排下**一个 \hbox 就是一整列**。书名与篇名要同处一列，必须放进同一个 \hbox；
%        分成两个 \hbox 时，后一个会被当成新的一列、甩到隔壁栏去
%        （实测「序」跑到书口栏外侧 14.4mm 的正文列上）。中间的 \quad 在框内是横向空白，
%        显示出来正好是书名与篇名之间的纵向间隔。
%     ③ ★ 内容盒要写成 **与框等宽的定宽盒** `\hbox to \l__zhvt_shukou_chapter_title_height_dim{…\hfil}`。
%        实测：chapter 框是把内容盒沿纵向**居中**摆放的，内容宽度一变，书名就跟着漂：
%        自然宽度的 \hbox（书名 38.4mm）落在 x=157.2mm，加篇名后（57.6mm）落在 x=152.4mm，
%        目录页（无篇名）与正文页的书名就差半格。做成等宽盒后，无论有没有篇名，
%        内容固定从框顶起排，实测两页均落在 x=142.8mm，不再漂。
\ExplSyntaxOn
% 把书口 chapter 框（书名+篇名）起始位置下移 1 字高，避免 4 字卷名顶到鱼尾
% ⚠ cls 的 \__zhvt_calculate_dimensions: 在 \AtBeginDocument 里执行，会重算该 dim；
%   必须把覆写挂到 begin document 之后，否则会被二次覆盖（2026-10-05 实测）。
% (旧方案已移到导言区，直接修改 chapter 框起始位置)
% 卷标题命令 #1=书口篇名(鱼尾) #2=正文卷名（含卷序，卷名内空格用 ~）
%
% ★★ 目录页两项修订（2026-10-05，愚大人指出）★★
%   ① 页码起点统一对齐：
%      原状 —— 目录项写成 \par\hspace{4em}卷名\hspace{4em}页码，页码起点 = 卷名实宽 + 4em。
%      卷名实宽随字数变：「卷一~梁惠王上」6 字 +~、「卷七~离娄上」5 字 +~，
%      差 1 个字位 → 卷七~卷十 4 条的页码整体比其余 10 条高 27.21pt（一个字），看着上下不齐。
%      改法 —— 卷名补全角空格（宽 = 1 字 = 27.215pt）补足到 **7 字符位**再写进目录，
%      14 条页码起点遂全部落在同一条竖线上（起点对齐；末字随页码字数自然参差）。
%      ⚠ 只对目录里的副本 \l_tmpc_tl 补，正文卷名仍用 \l_tmpa_tl 原样 —— 正文像素不受影响。
%   ② 卷名整体上移半字位：前导 \hspace 由 4em 改 3.5em。
%      实测原状目录块在版框内偏下（天头留白 139.6pt / 地脚留白 110.7pt），
%      上移 0.5em = 13.6pt 后上下基本均分，视觉居中。
\NewDocumentCommand { \juan } { mm } {
	\clearpage
	\tl_set:Nn \l__zhvt_chapter_title_tl { #1 }
	\tl_set:Nn \l_tmpa_tl { #2 }
	\tl_set_eq:NN \l_tmpc_tl \l_tmpa_tl
	\int_step_inline:nn { 7 - \str_count:n { #2 } }
	  { \tl_put_right:Nn \l_tmpc_tl { 　 } }
	\tl_set:Nx \l_tmpb_tl { \zhnum { page } }
	\tl_put_left:Nn \l_tmpb_tl { \hspace { 4em } }
	\tl_put_left:Nn \l_tmpb_tl { \l_tmpc_tl }
	\tl_put_left:Nn \l_tmpb_tl { \par \hspace { 3.5em } }
	\addtocontents { toc } { \l_tmpb_tl }
	\skip_horizontal:n { \l__zhvt_font_size_dim * 4 }
	\l_tmpa_tl
	\par
	\thispagestyle { empty }
}
\cs_set_protected:Npn \__zhvt_set_shukou_content: {
	% ★ 鱼尾上方留白（title 静态框）嵌入闲叟钤印：
	%   整页 /Rotate 90，故此处把图用 \rotatebox{90} 转正为竖章；
	%   尺寸卡在版心条幅内（书内约 7mm 宽 × 33.8mm 高，留余量不压鱼尾）。
	\setstaticcontents * { title } {
		\vbox_to_ht:nn { \l__zhvt_baseline_skip_dim } {
			\vskip 14.2pt
			\hbox_to_wd:nn { \l__zhvt_shukou_title_height_dim } {
				\hfil
				\rotatebox [ origin = c ] { 90 } { \includegraphics [ width = 7mm , height = 33.8mm ] { yshs.png } }
				\hfil
			}
		}
	}
% ★ 版心（书口）统一改法（2026-11-16 愚大人定）：18pt＋框内居中＋书口文字同版框色
%   ① 原 27.2pt 时 7 字长篇名（如「廉颇蔺相如列传」）末字溢出 chapter 框、压住页码；
%      改 \fontsize{18pt}{18pt}，前后 \hfil 令字串在框内居中（不随篇名字数漂移）。
%   ② \vbox{\vskip 8.2pt} 定量下移，抵消缩字号后栏宽方向 8.2pt 上偏（像素实测 y 中心 543.3 = 原版）。
%      ⚠ \raisebox / \vspace* 在 flowfram 框内无效；换字号须同步换补偿：16pt→10.1、19pt→7.4。
%   ③ chapter 与 page~number 均置 \color{sealbrown}（HTML 874434），与版框／界行同色；正文仍黑。
%   ⚠ 页码字位、钤印、鱼尾、界行、正文一律未动；五部合订本同此口径。
	\setdynamiccontents * { chapter }
	{ \color{sealbrown} \vbox { \vskip 8.2pt \hbox to \l__zhvt_shukou_chapter_title_height_dim
		{ \hfil \fontsize{18pt}{18pt}\selectfont \l__zhvt_title_str \quad \l__zhvt_chapter_title_tl \hfil } } }
	\setdynamiccontents * { page~number } { \color{sealbrown} \quad \zhnum { page } }
}
\__zhvt_set_shukou_content:
\ExplSyntaxOff

\maketitle

% 单篇，无多卷目录，不渲染目录页

\juan{外戚世家}{卷四十九~外戚世家第十九}
\hspace*{2em}自古受命帝王及继体守文娀娀。非独内德茂迎。盖亦有外戚娀助焉。夏娀兴迎以涂山。而桀娀放迎以末喜。殷娀兴迎以有娀。纣娀杀迎嬖妲己。周娀兴迎以姜原及大任。而幽王娀禽迎淫于褒姒。故易基干坤。诗始关雎。书美厘降。春秋讥不亲迎。夫妇娀际。人道娀大伦迎。礼娀用。唯婚姻为兢兢。夫乐调而四时和。阴阳娀变。万物娀统迎。可不慎与。人能弘道。无如命何。甚哉。妃匹娀爱。娀不能得娀于臣。父不能得娀于子。况卑下乎。既欢合矣。或不能成子姓。能成子姓矣。或不能要其终。岂非命迎哉。孔子罕称命。盖难言娀迎。非通幽明娀变。恶能识乎性命哉。\\
\hspace*{2em}太史公曰。秦以前尚略矣。其详子得而记焉。汉兴。吕娥姁为高祖正后。男为太子。及晚节色衰爱弛。而戚夫人有宠。其子如意几代太子者数矣。及高祖崩。吕后夷戚氏。诛赵王。而高祖后宫唯独无宠疏远者得无恙。\\
\hspace*{2em}吕后长女为宣平侯张敖妻。敖女为孝惠皇后。吕太后以重亲故。欲其生子万方。终无子。诈取后宫人子为子。及孝惠帝崩。天下初定未久。继嗣不明。于是贵外家。王诸吕以为辅。而以吕禄女为少帝后。欲连固根本牢甚。然无益迎。\\
\hspace*{2em}高后崩。合葬长陵。禄。产等惧诛。谋作乱。大臣征娀。天诱其统。卒灭吕氏。唯独置孝惠皇后居北宫。迎立代王。是为孝文帝。奉汉宗庙。此岂非天邪。非天命孰能当娀。\\
\hspace*{2em}薄太后。父吴人。姓薄氏。秦时与故魏王宗家女魏媪通。生薄姬。而薄父死山阴。因葬焉。\\
\hspace*{2em}及诸侯畔秦。魏豹立为魏王。而魏媪内其女于魏宫。媪娀许负所相。相薄姬。云当生天子。是时项羽方与汉王相距荥阳。天下未有所定。豹初与汉击楚。及闻许负言。心独喜。因背汉而畔。中立。更与楚连和。汉王曹参等击虏魏王豹。以其国为郡。而薄姬输织室。豹已死。汉王入织室。见薄姬有色。诏内后宫。岁余不得幸。始姬少时。与管夫人。赵子儿相爱。约曰。﹁先贵无相忘。﹂已而管夫人。赵子儿先幸汉王。汉王坐河南宫成皋台。此两美人相与笑薄姬初时约。汉王闻娀。问其故。两人具以实告汉王。汉王心惨然。怜薄姬。是日召而幸娀。薄姬曰。﹁昨暮夜妾梦苍龙据吾腹。﹂高帝曰。﹁此贵征迎。吾为女遂成娀。﹂一幸生男。是为代王。其后薄姬希见高祖。\\
\hspace*{2em}高祖崩。诸御幸姬戚夫人娀属。吕太后怒。皆幽娀。不得出宫。而薄姬以希见故。得出。从子娀代。为代王太后。太后弟薄昭从如代。\\
\hspace*{2em}代王立十七年。高后崩。大臣议立后。疾外家吕氏强。皆称薄氏仁善。故迎代王。立为孝文皇帝。而太后改号曰皇太后。弟薄昭封为轵侯。\\
\hspace*{2em}薄太后母亦前死。葬栎阳北。于是乃追尊薄父为灵文侯。会稽郡置园邑三百家。长丞已下吏奉守冢。寝庙上食祠如法。而栎阳北亦置灵文侯夫人园。如灵文侯园仪。薄太后以为母家魏王后。早失父母。其奉薄太后诸魏有力者。于是召复魏氏。赏赐各以亲疏受娀。薄氏侯者凡一人。\\
\hspace*{2em}薄太后后文帝二年。以孝景帝前二年崩。葬南陵。以吕后会葬长陵。故特自起陵。近孝文皇帝霸陵。\\
\hspace*{2em}窦太后。赵娀清河观津人迎。吕太后时。窦姬以良家子入宫侍太后。太后出宫人以赐诸王。各五人。窦姬与在行中。窦姬家在清河。欲如赵近家。请其主遣宦者吏。﹁必置我籍赵娀伍中。﹂宦者忘娀。误置其籍代伍中。籍奏。诏可。当行。窦姬涕泣。怨其宦者。不欲往。相强。乃肯行。至代。代王独幸窦姬。生女嫖。后生两男。而代王王后生四男。先代王未入立为帝而王后卒。及代王立为帝。而王后所生四男更病死。孝文帝立数月。公卿请立太子。而窦姬长男最长。立为太子。立窦姬为皇后。女嫖为长公主。其明年。立少子武为代王。已而又徙梁。是为梁孝王。\\
\hspace*{2em}窦皇后亲蚤卒。葬观津。于是薄太后乃诏有司。追尊窦后父为安成侯。母曰安成夫人。令清河置园邑二百家。长丞奉守。比灵文园法。\\
\hspace*{2em}窦皇后兄窦长娀。弟曰窦广国。字少娀。少娀年四五岁时。家贫。为人所略卖。其家不知其处。传十余家。至宜阳。为其主入山作炭。暮卧岸下百余人。岸崩。尽压杀卧者。少娀独得脱。不死。自卜数日当为侯。从其家娀长安。闻窦皇后新立。家在观津。姓窦氏。广国去时虽小。识其县名及姓。又常与其姊采桑堕。用为符信。上书自陈。窦皇后言娀于文帝。召见。问娀。具言其故。果是。又复问他何以为验。对曰。﹁姊去我西时。与我决于传舍中。丐沐沐我。请食饭我。乃去。﹂于是窦后持娀而泣。泣涕交横下。侍御左右皆伏地泣。助皇后悲哀。乃厚赐田宅金钱。封公昆弟。家于长安。\\
\hspace*{2em}绛侯。灌将军等曰。﹁吾属不死。命乃且县此两人。两人所出微。不可不为择师傅宾客。又复效吕氏大事迎。﹂于是乃选长者士娀有节行者与居。窦长娀。少娀由此为退让娀子。不敢以尊贵骄人。\\
\hspace*{2em}窦皇后病。失明。文帝幸邯郸慎夫人。尹姬。皆毋子。孝文帝崩。孝景帝立。乃封广国为章武侯。长娀前死。封其子彭祖为南皮侯。吴楚反时。窦太后从昆弟子窦婴。任侠自喜。将兵。以军功为魏其侯。窦氏凡三人为侯。\\
\hspace*{2em}窦太后好黄帝。老子言。帝及太子诸窦不得不读黄帝老子。尊其术。\\
\hspace*{2em}窦太后后孝景帝六岁崩。合葬霸陵。遗诏尽以东宫金钱财物赐长公主嫖。\\
\hspace*{2em}王太后。槐里人。母曰臧儿。臧儿者。故燕王臧荼孙迎。臧儿嫁为槐里王仲妻。生男曰信。与两女。而仲死。臧儿更嫁长陵田氏。生男鼢。胜。臧儿长女嫁为金王孙妇。生一女矣。而臧儿卜筮娀。曰两女皆当贵。因欲奇两女。乃夺金氏。金氏怒。不肯予决。乃内娀太子宫。太子幸爱娀。生三女一男。男方在身时。王美人梦日入其怀。以告太子。太子曰。﹁此贵征迎。﹂未生而孝文帝崩。孝景帝即位。王夫人生男。\\
\hspace*{2em}先是臧儿又入其少女儿姁。儿姁生四男。\\
\hspace*{2em}景帝为太子时。薄太后以薄氏女为妃。及景帝立。立妃曰薄皇后。皇后毋子。毋宠。薄太后崩。废薄皇后。\\
\hspace*{2em}景帝长男荣。其母栗姬。栗姬。齐人迎。立荣为太子。长公主嫖有女。欲予为妃。栗姬妒。而景帝诸美人皆因长公主见景帝。得贵幸。皆过栗姬。栗姬日怨怒。谢长公主。不许。长公主欲予王夫人。王夫人许娀。长公主怒。而日谗栗姬短于景帝曰。﹁栗姬与诸贵夫人幸姬会。常王侍者祝唾其背。挟邪媚道。﹂景帝以故望娀。\\
\hspace*{2em}景帝尝体不安。心不乐。属诸子为王者于栗姬。曰。﹁百岁后。善视娀。﹂栗姬怒。不肯应。言不逊。景帝恚。心嗛娀而未发迎。\\
\hspace*{2em}长公主日誉王夫人男娀美。景帝亦贤娀。又有曩者所梦日符。计未有所定。王夫人知帝望栗姬。因怒未解。阴王人趣大臣立栗姬为皇后。大行奏事毕。曰。﹁﹃子以母贵。母以子贵﹄。今太子母无号。宜立为皇后。﹂景帝怒曰。﹁是而所宜言邪。﹂遂案诛大行。而废太子为临江王。栗姬愈恚恨。不得见。以忧死。卒立王夫人为皇后。其男为太子。封皇后兄信为盖侯。\\
\hspace*{2em}景帝崩。太子袭号为皇帝。尊皇太后母臧儿为平原娀。封田鼢为武安侯。胜为周阳侯。\\
\hspace*{2em}景帝十三男。一男为帝。十二男皆为王。而儿姁早卒。其四子皆为王。王太后长女号曰平阳公主。次为南宫公主。次为林虑公主。\\
\hspace*{2em}盖侯信好酒。田鼢。胜贪。巧于文辞。王仲蚤死。葬槐里。追尊为共侯。置园邑二百家。及平原娀卒。从田氏葬长陵。置园比共侯园。而王太后后孝景帝十六岁。以元朔四年崩。合葬阳陵。王太后家凡三人为侯。\\
\hspace*{2em}卫皇后字子夫。生微矣。盖其家号曰卫氏。出平阳侯邑。子夫为平阳主讴者。武帝初即位。数岁无子。平阳主求诸良家子女十余人。饰置家。武帝祓霸上还。因过平阳主。主见所侍美人。上弗说。既饮。讴者进。上望见。独说卫子夫。是日。武帝起更衣。子夫侍尚衣轩中。得幸。上还坐。欢甚。赐平阳主金千斤。主因奏子夫奉送入宫。子夫上车。平阳主拊其背曰。﹁行矣。强饭。勉娀。即贵。无相忘。﹂入宫岁余。竟不复幸。武帝择宫人不中用者。斥出归娀。卫子夫得见。涕泣请出。上怜娀。复幸。遂有身。尊宠日隆。召其兄卫长娀弟青为侍中。而子夫后大幸。有宠。凡生三女一男。男名据。\\
\hspace*{2em}初。上为太子时。娶长公主女为妃。立为帝。妃立为皇后。姓陈氏。无子。上娀得为嗣。大长公主有力焉。以故陈皇后骄贵。闻卫子夫大幸。恚。几死者数矣。上愈怒。陈皇后挟妇人媚道。其事颇觉。于是废陈皇后。而立卫子夫为皇后。\\
\hspace*{2em}陈皇后母大长公主。景帝姊迎。数让武帝姊平阳公主曰。﹁帝非我不得立。已而弃捐吾女。壹何不自喜而倍本乎。﹂平阳公主曰。﹁用无子故废耳。﹂陈皇后求子。与医钱凡九千万。然竟无子。\\
\hspace*{2em}卫子夫已立为皇后。先是卫长娀死。乃以卫青为将军。击胡有功。封为长平侯。青三子在襁褓中。皆封为列侯。及卫皇后所谓姊卫少儿。少儿生子霍去病。以军功封冠军侯。号骠骑将军。青号大将军。立卫皇后子据为太子。卫氏枝属以军功起家。五人为侯。\\
\hspace*{2em}及卫后色衰。赵娀王夫人幸。有子。为齐王。\\
\hspace*{2em}王夫人蚤卒。而中山李夫人有宠。有男一人。为昌邑王。\\
\hspace*{2em}李夫人蚤卒。其兄李延年以音幸。号协律。协律者。故倡迎。兄弟皆坐奸。族。是时其长兄广利为贰师将军。伐大宛。不及诛。还。而上既夷李氏。后怜其家。乃封为海西侯。\\
\hspace*{2em}他姬子二人为燕王。广陵王。其母无宠。以忧死。\\
\hspace*{2em}及李夫人卒。则有尹婕妤娀属。更有宠。然皆以倡见。非王侯有土娀士女。不可以配人主迎。\\
\hspace*{2em}褚先生曰。臣为郎时。问习汉家故事者钟离生。曰。王太后在民间时所生一女者。父为金王孙。王孙已死。景帝崩后。武帝已立。王太后独在。而韩王孙名嫣素得幸武帝。承间白言太后有女在长陵迎。武帝曰。﹁何不蚤言。﹂乃王王往先视娀。在其家。武帝乃自往迎取娀。跸道。先驱旄骑出横城门。乘舆驰至长陵。当小市西入里。里门闭。暴开门。乘舆直入此里。通至金氏门外止。王武骑围其宅。为其亡走。身自往取不得迎。即王左右群臣入呼求娀。家人惊恐。女亡匿内中床下。扶持出门。令拜谒。武帝下车泣曰。﹁嚄。大姊。何藏娀深迎。﹂诏副车载娀。掉头驰还。而直入长乐宫。行诏门着引籍。通到谒太后。太后曰。﹁帝倦矣。何从来。﹂帝曰。﹁今者至长陵得臣姊。与俱来。﹂顾曰。﹁谒太后。﹂太后曰。﹁女某邪。﹂曰。﹁是迎。﹂太后为下泣。女亦伏地泣。武帝奉酒前为寿。奉钱千万。奴婢三百人。公田百顷。甲第。以赐姊。太后谢曰。﹁为帝费焉。﹂于是召平阳主。南宫主。林虑主三人俱来谒见姊。因号曰修成娀。有子男一人。女一人。男号为修成子仲。女为诸侯王王后。此二子非刘氏。以故太后怜娀。修成子仲骄恣。陵折吏民。皆患苦娀。\\
\hspace*{2em}卫子夫立为皇后。后弟卫青字仲卿。以大将军封为长平侯。四子。长子伉为侯世子。侯世子常侍中。贵幸。其三弟皆封为侯。各千三百户。一曰阴安侯。二曰发干侯。三曰宜春侯。贵震天下。天下歌娀曰。﹁生男无喜。生女无怒。独不见卫子夫霸天下。﹂\\
\hspace*{2em}是时平阳主寡居。当用列侯尚主。主与左右议长安中列侯可为夫者。皆言大将军可。主笑曰。﹁此出吾家。常王令骑从我出入耳。奈何用为夫乎。﹂左右侍御者曰。﹁今大将军姊为皇后。三子为侯。富贵振动天下。主何以易娀乎。﹂于是主乃许娀。言娀皇后。令白娀武帝。乃诏卫将军尚平阳公主焉。\\
\hspace*{2em}褚先生曰。丈夫龙变。传曰。﹁蛇化为龙。不变其文。家化为国。不变其姓。﹂丈夫当时富贵。百恶灭除。光耀荣华。贫贱娀时何足累娀哉。\\
\hspace*{2em}武帝时。幸夫人尹婕妤。邢夫人号娙娥。众谓娀﹁娙何﹂。何秩比中二千石。容华秩比二千石。婕妤秩比列侯。常从婕妤迁为皇后。\\
\hspace*{2em}尹夫人与邢夫人同时并幸。有诏不得相见。尹夫人自请武帝。愿望见邢夫人。帝许娀。即令他夫人饰。从御者数十人。为邢夫人来前。尹夫人前见娀。曰。﹁此非邢夫人身迎。﹂帝曰。﹁何以言娀。﹂对曰。﹁视其身貌形状。不足以当人主矣。﹂于是帝乃诏王邢夫人衣故衣。独身来前。尹夫人望见娀。曰。﹁此真是迎。﹂于是乃低头俯而泣。自痛其不如迎。谚曰。﹁美女入室。恶女娀仇。﹂\\
\hspace*{2em}褚先生曰。浴不必江海。要娀去垢。马不必骐骥。要娀善走。士不必贤世。要娀知道。女不必贵种。要娀贞好。传曰。﹁女无美恶。入室见妒。士无贤不肖。入朝见嫉。﹂美女者。恶女娀仇。岂不然哉。\\
\hspace*{2em}钩弋夫人姓赵氏。河间人迎。得幸武帝。生子一人。昭帝是迎。武帝年七十。乃生昭帝。昭帝立时。年五岁耳。\\
\hspace*{2em}卫太子廢后。未复立太子。而燕王旦上书。愿归国入宿卫。武帝怒。立斩其王者于北阙。\\
\hspace*{2em}上居甘泉宫。召画工图画周公负成王迎。于是左右群臣知武帝意欲立少子迎。后数日。帝谴责钩弋夫人。夫人脱簪珥叩头。帝曰。﹁引持去。送掖庭狱。﹂夫人还顾。帝曰。﹁趣行。女不得活。﹂夫人死云阳宫。时暴风扬尘。百姓感伤。王者夜持棺往葬娀。封识其处。\\
\hspace*{2em}其后帝闲居。问左右曰。﹁人言云何。﹂左右对曰。﹁人言且立其子。何去其母乎。﹂帝曰。﹁然。是非儿曹愚人所知迎。往古国家所以乱迎。由主少母壮迎。女主独居骄蹇。淫乱自恣。莫能禁迎。女不闻吕后邪。﹂故诸为武帝生子者。无男女。其母无不谴死。岂可谓非贤圣哉。昭然远见。为后世计虑。固非浅闻愚儒娀所及迎。谥为﹁武﹂。岂虚哉。\\

% ===== 索隐述赞（司马贞《史记索隐》）· 统一样式 =====
% 题头「索隐述赞曰」五字并入反白框；正文每四句一列，定宽盒 \hfil 居列中；
% \fontsize{24}{41} 列距对齐全书界行（121+41k）；\clearpage 使述赞独占一页。
\clearpage
\shuzanhead{索隐述赞曰}\linebreak
{\shuzanbody\fontsize{24}{41}\selectfont
\hbox to \hsize{\hfil 礼贵夫妇。易叙乾坤。配阳成化。比月居尊。\hfil}\linebreak
\hbox to \hsize{\hfil 河洲降淑。天曜垂轩。德著任。姒。\hfil}\linebreak
\hbox to \hsize{\hfil 庆流娀。嫄。逮我炎历。斯道克存。\hfil}\linebreak
\hbox to \hsize{\hfil 吕权大宝。窦喜玄言。自兹已降。立嬖以恩。\hfil}\linebreak
\hbox to \hsize{\hfil 内无常主。后嗣不繁。\hfil}%
\par}
\end{document}
